\section{Численное интегрирование}
\label{sec:numerical-integration}

\subsection{Квадратурные формулы}

Численное интегрирование заменяет интеграл конечной суммой:
\[
    I(f)=\int_a^b w(x)f(x)\,dx
    \approx
    Q_n(f)=\sum_{i=1}^{n}A_i f(x_i),
\]
где $x_i$ --- узлы, $A_i$ --- веса, $w(x)$ --- весовая функция.
Погрешность равна
\[
    R_n(f)=I(f)-Q_n(f).
\]
Алгебраическая степень точности квадратурной формулы --- наибольшая степень
$p$, для которой формула точна для всех многочленов степени не выше $p$.

\subsection{Формулы Ньютона--Котеса}

Они получаются интегрированием интерполяционного многочлена на равноотстоящих
узлах. Простейшие формулы на одном отрезке:

Левые прямоугольники:
\[
    \int_a^b f(x)\,dx
    \approx (b-a)f(a),
    \qquad R=O((b-a)^2).
\]

Средние прямоугольники:
\[
    \int_a^b f(x)\,dx
    \approx (b-a)f\left(\frac{a+b}{2}\right),
\]
\[
    R=\frac{(b-a)^3}{24}f''(\xi).
\]

Трапеции:
\[
    \int_a^b f(x)\,dx
    \approx \frac{b-a}{2}\bigl(f(a)+f(b)\bigr),
\]
\[
    R=-\frac{(b-a)^3}{12}f''(\xi).
\]

Формула Симпсона:
\[
    \int_a^b f(x)\,dx
    \approx
    \frac{b-a}{6}
    \left[f(a)+4f\left(\frac{a+b}{2}\right)+f(b)\right],
\]
\[
    R=-\frac{(b-a)^5}{2880}f^{(4)}(\xi).
\]
Симпсон точен для многочленов степени до трёх включительно.

\subsection{Составные формулы}

Разобьём $[a,b]$ на $N$ равных частей, $h=(b-a)/N$.
Составная формула трапеций:
\[
    T_h=h\left[
    \frac{f_0+f_N}{2}+\sum_{i=1}^{N-1}f_i
    \right],
\]
\[
    I-T_h=-\frac{b-a}{12}h^2f''(\xi).
\]

Для чётного $N$ составная формула Симпсона:
\[
S_h=\frac{h}{3}
\left[
 f_0+f_N
 +4\sum_{\substack{i=1\\i\text{ нечётное}}}^{N-1}f_i
 +2\sum_{\substack{i=2\\i\text{ чётное}}}^{N-2}f_i
\right],
\]
\[
    I-S_h=-\frac{b-a}{180}h^4f^{(4)}(\xi).
\]

\subsection{Правило Рунге и адаптивное интегрирование}

Если погрешность метода имеет вид $Ch^p$, то по двум сеткам:
\[
    I-Q_{h/2}\approx\frac{Q_{h/2}-Q_h}{2^p-1}.
\]
Это правило Рунге используется для оценки ошибки. В адаптивном алгоритме отрезок
делят только там, где локальная оценка превышает допуск. Адаптивный Симпсон
сравнивает формулу на всём отрезке с суммой формул на двух половинах.

\subsection{Квадратуры Гаусса}

В формулах Ньютона--Котеса узлы фиксированы. Если выбирать и узлы, и веса, то
$n$-узловая квадратура может быть точна для многочленов степени до $2n-1$:
\[
    \int_a^b w(x)f(x)\,dx
    \approx \sum_{i=1}^{n}A_i f(x_i).
\]

\textbf{Теорема Гаусса--Кристоффеля.}
Пусть вес $w(x)$ положителен почти всюду на $[a,b]$ и существуют необходимые
моменты. Тогда узлы $n$-точечной квадратуры Гаусса наивысшей алгебраической
степени точности $2n-1$ являются корнями ортогонального многочлена степени $n$
относительно веса $w(x)$ на $[a,b]$, а квадратурные веса $A_i$ положительны.

Для $w(x)=1$ на $[-1,1]$ узлами являются корни многочлена Лежандра $P_n$.
Переход на произвольный отрезок:
\[
    x=\frac{a+b}{2}+\frac{b-a}{2}t.
\]

Пример двухузловой формулы Гаусса--Лежандра:
\[
    \int_{-1}^{1}f(x)\,dx
    \approx
    f\left(-\frac1{\sqrt3}\right)
    +f\left(\frac1{\sqrt3}\right).
\]
Она точна для многочленов степени до трёх.

\subsection{Специальные квадратуры}

Различным весам и областям соответствуют семейства ортогональных многочленов:
\begin{itemize}
    \item Гаусс--Чебышёв:
    \[
        \int_{-1}^{1}\frac{f(x)}{\sqrt{1-x^2}}\,dx;
    \]
    \item Гаусс--Лагерр:
    \[
        \int_{0}^{\infty}e^{-x}f(x)\,dx;
    \]
    \item Гаусс--Эрмит:
    \[
        \int_{-\infty}^{\infty}e^{-x^2}f(x)\,dx.
    \]
\end{itemize}
Выбор квадратуры, согласованной с весом и особенностями подынтегральной функции,
резко повышает эффективность.

\subsection{Несобственные интегралы и особенности}

Для бесконечного интервала применяют замену переменной или специальные
квадратуры Лагерра и Эрмита. При интегрируемой особенности полезно выделить её
аналитически или сделать замену, сглаживающую функцию. Простое равномерное
сгущение сетки часто неэффективно.

Для периодических гладких функций составная формула трапеций может сходиться
очень быстро. Для осциллирующих интегралов применяют методы Филона и квадратуры,
учитывающие фазу.

\subsection{Ошибки вычислений}

Численное интегрирование обычно сглаживает шум, однако при большом числе узлов
накапливаются ошибки округления. Формулы Ньютона--Котеса высокого порядка могут
иметь большие веса разных знаков и быть неустойчивыми. Поэтому на практике чаще
используют составные низкопорядковые формулы, адаптивные методы или квадратуры
Гаусса.

\subsection{Что важно сказать на экзамене}

Нужно привести общий вид квадратурной формулы, формулы трапеций и Симпсона с
порядками ошибок, правило Рунге и принцип квадратур Гаусса. Главное отличие:
Ньютон--Котес использует равноотстоящие узлы, Гаусс выбирает оптимальные узлы и
достигает точности $2n-1$.

\newpage